
This shall be a 1 page-long comprehensive introduction to the topic of this thesis. The introduction should state the research/practical goals of this thesis.

Comment: Research goals roughly, chatgpt. I ahave dumped all my thoughts in one go and I have given you context here

{
This thesis started with examining gnn based methods in reconstruction for ldmx, especially for the ecal and trigger which the supervisor has better understanding of. The trigger, reconstruction methods are key for the physics reach of the experiment. I wanted to research ml-based methods because I am a ml engineer but that is not a good reason for research but that is merely practical context for you chatgpt. 

It started out here with gnn because they have a very natural interpretation for hits in the ecal and for detector hits and detector derived objects from heterogenous amount of detectors. These come as unorded sets as input and the gnn has no inherent structure bias, as compared to the cnn for example. This research interest also aligns with the filed of hep where gnn have had an interest to study. However when I started with this and I tried to read up on the topic I quickly discovered mlpf, a ml based particleflow method at CSM CERN and I identified a very very similar research goal of finding ways to utilize ml and the strengths of ml in the context of hep, but of course for a completely different experiment, detector and for potentially different assignments for the method as well... 

This led to that I wanted to see ways of doing hit-level analysis (actually called node classification for shower separation and node regression for other regression based tasks), since that had not been done before by previous students who have studied ways of using ml in ldmx. Previous students have mostly focused on electron counting with ml based methods, which are event level, but I wanted to see ways of doing more in depth analysis being hit level. I also identified that there is a recurring interest among students (and an ongoing interest among researchers and ldmx collaborators) of using ML and in the same manner that GNN are interesting because they are flexibile in their input, they are also kind of a flexible tool where you should be able to re-use a platform that has implemented a solution of how to input unorded detector sets where you can test new hypotheses, test inserting new information to your models, test assigning new tasks to your model or even try to design a coherent multi-task model performing full event particleflow reconstruction similar that the mlpf does at cms cern! So there are several research goals and also a practical goal where I aim to build a platform that can be re-used (which is big!). Being able to make hit level node classification will take some coding leg work for sure. This was also motivated by the fact that I saw the previous students work but I could not really re-use their code again I dont want that to happen again! I also recognize that my work will probably take a lot of time compared to my scope of merely 30 credits since I would have loved to make a full event reconstructor like mlpf cms cern but I do not have the time to reach there I have to work in steps and thus limit my study to the first logical step of hit-level node-classification which is arguably probably shower separation of 2e and 3e pileup events (pileup events are perhaps where ml models are the most interesting and potentially the most effective anyways). 

Research goals include thus how effective ml models are at shower separation. If we trust that the models have found an optimal solution, does this study then find the boundary of how much you can reconstruct given an event? This boundary is interesting since it sets the boundary of which events you want to trigger and it sets the boundary of how much analysis you can recover from events restricting your physics reach given limited amount of time and resources for an experiment. You could try to achieve this boundary with rule-based reconstruction but if this boundary of ambiguity that you simply cannot break then you will not be able to breach that with rule based methods either this study could inform a task of optimizing rule-based reconstruction if that is desired since we eliminate the black box nature. A research goal here is also to study if/how the trigger scintillator subdetector can improve analysis in the ecal. If it does not improve, what does that mean? If we trust that the models are optimal, then ts cannot help, but if we dont trust that we can expand the study and try to include object level information forcing the models to find new ways of improving shower separation. If we fastly build a model that counts the number of electrons, then does the ts help with that? A research goal is to find what factors that restrict the models ability to separate showers and once again if we trust that the models are optimal this gives us a boundary of where ambiguity cannot be breached. 

A practical/research goal is to map and find what the next steps would be for looking for ways to utilize ml based gnn/transformer based methods and how these seem useful or not.

One research goal is to compare the gnn and the transformer architecture and try to evaluate which architecture seem to be the best/most suitable to use? Compare their performances to see if one is better suited for this task in the sense that it performs better, and if they perform the same, which is preferred in that case and are there other considerations that make one or the other more preferrable? 

}


\textbf{Content, in essence:}

This thesis investigates modern machine learning architectures including GNN-based and transformer-based reconstruction methods in particle physics for the Light Dark Matter eXperiment (LDMX), specifically tested under simulated pile-up conditions. The work is motivated by the search for light dark matter  and by recent progress in machine-learning-based event reconstruction in high-energy physics, particularly approaches inspired by machine-learned particle flow.


\textbf{The introduction should include an explanation for what reconstruction is in HEP and relevant for my thesis of ML based reconstruction and using heterogenrous subdetector data. This way I dont have to define or explain the term in a later stage and so far I more or less assume that the reader knows what reconstruction is. }



